\section{How Updates Move}
\label{sec:updates}

This section describes the algebra each contract carries, the synchronization protocol that propagates it efficiently across replicas, the co-host mesh through which updates flow, and the leases that keep hosting aligned with demand.

\subsection{Contract State as an Idempotent Commutative Monoid}

\begin{definition}[Contract state algebra]
\label{def:algebra}
A contract $\contract$ defines:
\begin{enumerate}[leftmargin=1.4em, itemsep=2pt, topsep=2pt]
  \item a \emph{state space} $\statespace_\contract$ (an abstract set; the wire encoding is a byte string and the contract code defines the encoding);
  \item a \emph{validity predicate} $\valid_\contract : \statespace_\contract \to \{\bot, \top\}$;
  \item an \emph{identity element} $e_\contract \in \statespace_\contract$ with $\valid_\contract(e_\contract) = \top$;
  \item a binary \emph{merge} operation $\merge_\contract : \statespace_\contract \times \statespace_\contract \to \statespace_\contract$, which is associative, commutative, and idempotent, with $e_\contract$ as identity.
\end{enumerate}
\end{definition}

An idempotent commutative monoid is equivalently a join-semilattice; the partial order is $a \sqsubseteq b \iff a \merge_\contract b = b$. This is the algebraic structure of a state-based CRDT \citep{shapiro2011crdt}. The monoid presentation emphasizes the relevant operational property: peers can merge state at any time, in any order, with any multiplicity, and obtain the same result.

The validity predicate is checked at the boundary. A peer that receives a candidate state $\sigma$ with $\valid_\contract(\sigma) = \bot$ rejects it and never merges it into its local replica. Validity typically encodes integrity invariants beyond what the type system enforces, for example ``every record in the log carries a valid signature by an authorized key.''

An \emph{update} is itself an element of $\statespace_\contract$. When a peer holding state $\sigma$ receives an update $u$, it computes $\sigma' = \sigma \merge_\contract u$, subject to validity. There is no separate notion of operation: the algebra of state is the algebra of change.

\begin{remark}[Convergence]
\label{rem:convergence}
For a finite multiset of valid updates $u_1, \ldots, u_n \in \statespace_\contract$ and a starting state $\sigma_0$, the merged result $\sigma_0 \merge_\contract u_{\pi(1)} \merge_\contract \cdots \merge_\contract u_{\pi(n)}$ is independent of the permutation $\pi$ by associativity and commutativity, and idempotence makes duplicate delivery harmless. These are the sufficient conditions for the Strong Eventual Consistency theorem of Shapiro et al.\ \citep{shapiro2011crdt}: peers that have delivered the same set of valid updates converge regardless of order or repetition. Freenet applies this result to monoids supplied by individual contracts.
\end{remark}

\subsection{Worked Examples}

\paragraph{Monotone counter.} $\statespace = \mathbb{N}$, $\merge = \max$, $e = 0$. The partial order is the usual $\leq$.

\paragraph{Last-writer-wins register.} $\statespace = V \times \mathbb{N} \times H$, where $V$ is the value type, $\mathbb{N}$ a Lamport-style timestamp, and $H$ a tiebreaker hash. Merge is lexicographic max on the (timestamp, hash) pair.

\paragraph{Signed log.} $\statespace$ is a finite set of records, each record carrying a signature by an authorized key. $\merge = \cup$, $e = \emptyset$. Validity rejects any state containing a record whose signature does not verify against a key listed in $\mathrm{params}$.

\paragraph{Observed-remove set.} $\statespace = (E, T)$ where $E$ is a set of $(\mathrm{elem}, \mathrm{id})$ pairs and $T$ a set of $\mathrm{id}$s that have been tombstoned. $\merge$ is pairwise union; the materialized set is $\{\mathrm{elem} : (\mathrm{elem}, \mathrm{id}) \in E,\ \mathrm{id} \notin T\}$.

A bounded signed log (for example, one that keeps the most recent $N$ records) requires an eviction function determined by the merged set. If eviction instead depends on arrival order, the merge ceases to be idempotent. Contract authors are responsible for maintaining this property.

\paragraph{Discarding requires a permanent condition.} The eviction rule above generalizes. Whenever a contract discards part of a merged result --- an entry that fails a predicate, a record whose author is not authorized --- the condition it discards on must be one that further growth cannot reverse. Validity may depend on the rest of the state, so a component invalid at $\sigma$ can become valid at some $\sigma' \sqsupseteq \sigma$; discarding it at $\sigma$ makes the result depend on the order in which the two arrived, and nothing revisits the decision later. The observed-remove set above is the sound shape: a tombstone is permanent, so materializing $E \setminus T$ discards on a condition that cannot flip. A record dropped because its author has not yet been merged is the unsound shape, while the same rule applied to an author who has been tombstoned is sound --- the two are distinguished by whether the absence is final, not by whether the author is currently present. Where the condition is transient, the component must be retained and its visibility derived from the current state, or the update rejected in full so that its sender re-offers it.

\subsection{Where This Sits in the Design Space}

Bayou \citep{terry1995bayou} introduced application-supplied merge procedures in 1995, and CouchDB, Riak, and Git's merge drivers also expose conflict resolution to application code. Freenet combines this approach with constraints required by an untrusted peer-to-peer environment.

Contracts bind merge code to the data identity through content addressing, execute it in a sandbox with fuel and memory bounds, and require it to satisfy the idempotent-commutative-monoid laws. Bayou permits more general merge procedures, while Freenet uses these restrictions to support convergence on untrusted peers.

Fixed-lattice CRDT frameworks such as Automerge \citep{automerge}, Y.js \citep{yjs}, and Antidote \citep{akkoorath2016cure} compose applications from a predefined catalog. Freenet allows a chat room, reputation system, or public file feed to define a lattice specific to its consistency requirements. The runtime consequently cannot verify the algebraic properties statically. A violation appears as failure to converge rather than as a structural error at deployment time.

Freenet therefore permits more flexibility than fixed-lattice frameworks and less than Bayou. The summary/delta interface developed next also lets contracts specialize their synchronization protocols.

\subsection{Summary/Delta Synchronization}

Pairwise merge is sufficient for convergence but, on its own, would require peers to ship full state when reconciling. For non-trivial contracts this is prohibitive. The platform therefore extends the contract interface with three additional functions that allow peers to exchange only the difference between their states.

\begin{definition}[Sync interface]
\label{def:sync}
In addition to the algebra of \cref{def:algebra}, a contract provides
\begin{align*}
  \summarize_\contract &: \statespace_\contract \to \Sigma_\contract,\\
  \getdelta_\contract &: \statespace_\contract \times \Sigma_\contract \to \Delta_\contract,\\
  \applydelta_\contract &: \statespace_\contract \times \Delta_\contract \to \statespace_\contract,
\end{align*}
where $\Sigma_\contract$ and $\Delta_\contract$ are contract-defined byte-string types.
\end{definition}

The summary type $\Sigma_\contract$ is intended to be much smaller than full state. For a signed log it might be the set of record identifiers, or a Merkle root over them; for a counter it might be the counter's value; for a large observed-remove set it might be a Bloom filter, or a pair of Bloom filters over $E$ and $T$.

For the two-step protocol below to be useful, the contract author is required to ensure that $\getdelta$ produces something that, when applied, dominates the other peer's state in the partial order:

\begin{property}[Sync soundness]
\label{prop:sync}
For all $\sigma_A, \sigma_B \in \statespace_\contract$, let $s = \summarize_\contract(\sigma_B)$ and $\delta = \getdelta_\contract(\sigma_A, s)$. Then
\[
  \applydelta_\contract(\sigma_B, \delta) \;\merge_\contract\; \sigma_A
  \;=\;
  \applydelta_\contract(\sigma_B, \delta).
\]
\end{property}

\noindent Equivalently, applying $\delta$ at $\sigma_B$ brings $\sigma_B$ at least as far up the lattice as $\sigma_A$ would have if merged in full. The platform does not verify this property; the contract is responsible for it.

\paragraph{Deltas outlive their summaries.} \cref{prop:sync} fixes $\sigma_B$, but $s$ is a summary that $B$ published earlier, and $\delta$ may be applied after $B$ has merged further updates. The requirement is therefore the stronger one obtained by quantifying over every state above $\sigma_B$: for all $\sigma_B' \sqsupseteq \sigma_B$,
\[
  \applydelta_\contract(\sigma_B', \delta) \;\merge_\contract\; \sigma_A
  \;=\;
  \applydelta_\contract(\sigma_B', \delta).
\]
A contract meets this by computing $\delta$ from what $s$ shows the recipient to lack: the summary is what keeps $\delta$ small, and the lattice order is what keeps it applicable once the recipient has advanced. A delta whose effect depends on the recipient being exactly at $\sigma_B$ is a patch rather than a delta in this sense, and out-of-order delivery will misapply it.

\subsection{The Two-Step Sync Protocol}

% Two-step summary/delta sync protocol between peers A and B. Sequence
% diagram. A sends summarize(sigma_A); B computes delta from its state
% relative to A's summary; B sends delta back; A applies delta.
\begin{figure}[t]
\centering
\begin{tikzpicture}[
    peer/.style={draw=black!75, rounded corners=2pt, fill=blue!8,
                 minimum width=1.6cm, minimum height=0.7cm,
                 font=\small\bfseries, align=center},
    state/.style={font=\scriptsize, gray!35!black, align=center},
    msg/.style={->, >=Stealth, thick, black!75},
    msglabel/.style={font=\scriptsize, fill=white, inner sep=2pt, align=center},
    lifeline/.style={black!45, dashed},
    >=Stealth,
]
  % Peers A and B at top of lifelines.
  \node[peer] (A) at (0,0) {peer $A$};
  \node[peer] (B) at (8,0) {peer $B$};

  % Initial states (placed just below each peer box, outside it).
  \node[state, anchor=north] at (A.south) [yshift=-3pt] {state $\sigma_A$};
  \node[state, anchor=north] at (B.south) [yshift=-3pt] {state $\sigma_B$};

  % Lifelines drop down from below each peer box.
  \draw[lifeline] ([yshift=-22pt] A.south) -- ++(0,-4.4);
  \draw[lifeline] ([yshift=-22pt] B.south) -- ++(0,-4.4);

  % Step 1: A computes summary, sends to B.
  \node[state, anchor=east, align=right]
       at ($(A)+(-0.15,-1.0)$) {$s_A \gets \summarize_\contract(\sigma_A)$};
  \draw[msg] ($(A)+(0,-1.4)$) -- ($(B)+(0,-1.4)$)
       node[msglabel, midway] {$s_A$};

  % Step 2: B computes delta and sends back.
  \node[state, anchor=west, align=left]
       at ($(B)+(0.15,-2.2)$) {$\delta_B \gets \getdelta_\contract(\sigma_B,\, s_A)$};
  \draw[msg] ($(B)+(0,-2.7)$) -- ($(A)+(0,-2.7)$)
       node[msglabel, midway] {$\delta_B$};

  % Step 3: A applies delta.
  \node[state, anchor=east, align=right]
       at ($(A)+(-0.15,-3.4)$) {$\sigma_A \gets \applydelta_\contract(\sigma_A,\, \delta_B)$};

  % Step 4: (optional) symmetric reverse direction.
  \draw[msg, dashed] ($(B)+(0,-4.1)$) -- ($(A)+(0,-4.1)$)
       node[msglabel, midway, text=gray!45!black] {$s_B$};
  \draw[msg, dashed] ($(A)+(0,-4.6)$) -- ($(B)+(0,-4.6)$)
       node[msglabel, midway, text=gray!45!black] {$\delta_A$};
  \node[state, anchor=west, gray!45!black, align=left]
       at ($(B)+(0.15,-4.85)$) {(symmetric:\ $B$ converges\\toward $A$ in the same way)};

  % Per-round bandwidth annotation.
  \node[font=\scriptsize, gray!35!black, anchor=north, align=center]
       at ($(A)!0.5!(B)+(0,-5.4)$)
       {per-round bandwidth bounded by $|s_A| + |s_B| + |\delta_A| + |\delta_B|$,
        each application-defined};

\end{tikzpicture}
\caption{The two-step summary/delta sync protocol. Peer $A$ sends a
small summary $s_A$ of its state. Peer $B$ uses $s_A$ to compute a
\emph{delta} relative to its own state and ships only the difference
back, which $A$ applies under the contract's merge. The exchange is
symmetric (dashed); each direction runs over the same transport. The
contract supplies $\summarize_\contract$, $\getdelta_\contract$, and
$\applydelta_\contract$, so the protocol is generic over whatever
state representation the contract uses.}
\label{fig:summary-delta}
\end{figure}


Two peers $A$ and $B$ reconcile their states for a contract $\contract$ as follows:

\begin{enumerate}[leftmargin=1.4em, itemsep=2pt]
  \item $A$ computes $s_A = \summarize_\contract(\sigma_A)$ and sends $s_A$ to $B$.
  \item $B$ computes $\delta_B = \getdelta_\contract(\sigma_B, s_A)$ and sends $\delta_B$ to $A$. $A$ applies the delta: $\sigma_A \gets \applydelta_\contract(\sigma_A, \delta_B)$.
\end{enumerate}

The protocol is symmetric: $B$ runs the same exchange in the other direction. In practice the two directions are interleaved over a single connection. Bandwidth per reconciliation round is bounded by $|s_A| + |s_B| + |\delta_A| + |\delta_B|$, each of which is application-defined. Contracts that use approximate summaries may require multiple rounds before replicas fully converge (see the next subsection), in which case the total bandwidth is the sum of these per-round bounds.

\subsection{Malformed Summaries and Malicious Deltas}

The basic protocol reconciles honest replicas with stale state. Approximate summaries and malicious peers require additional consideration.

\paragraph{Approximate summaries.} A summary type such as a Bloom filter has false positives: $B$ may believe $A$ already has a record that in fact $A$ does not. In the bare protocol, the record is then not included in the delta and the two replicas remain divergent. Contracts that use approximate summaries must therefore either fall back to a more precise second round when divergence is suspected, or include the false-positive rate in their convergence guarantees and rely on subsequent reconciliation rounds (driven by later updates, subscription renewals, or explicit re-sync) to repair divergence. The platform offers no built-in fallback; the contract designs the round structure. Diagnosing a faulty or insufficient round structure is one of the open problems noted in \cref{sec:status}.

\paragraph{Malicious peers.} A peer that supplies a deliberately malformed summary, or returns a delta that violates the validity predicate, is treated like any other source of invalid state. The receiver applies $\valid_\contract$ to the merged result; if the result fails validity, the delta is rejected. A more subtle attack is for a peer to supply a delta that is valid but withholding (claiming to be up to date when it is not). The protocol cannot distinguish such a peer from one that has simply not yet received the missing updates; the replication factor across multiple peers near $\ell(\contract)$ is the structural defense, not the sync protocol itself.

\subsection{Subscriptions and the Co-Host Mesh}

Once a contract is published, each hosting peer advertises that fact to its connected neighbors. An update arriving at any host is merged into local state and fanned out to neighbors that also advertise themselves as hosts of the contract, with each pairwise exchange using the summary/delta protocol above. These advertised co-host links form the update mesh. The live fan-out target set is derived from co-host advertisements rather than from a separate subscriber forwarding list; periodic reconciliation of the advertisements repairs a missed announcement or retraction.

Subscriptions serve a different purpose: they express demand and maintain soft-state paths from interested peers toward the contract's home location. A peer on such a path hosts and advertises the contract, so the path becomes part of the same co-host mesh as every other replica. Hosting follows demand and eviction, while subscriptions are explicit and lease-bound.

Because $\merge_\contract$ is commutative, simultaneous updates injected at multiple points of the mesh do not require coordination. They propagate across co-host links and merge; the order in which different updates arrive at a given node does not affect its final state. In the deployed configuration a subscription lease lasts 8 minutes and a background task renews it every 2 minutes --- but only while the contract is still \emph{in use} at that peer, meaning it has a local client subscription, a registered downstream subscriber, or recent genuine \textsf{GET} or \textsf{PUT} access. A hosting peer with no remaining interest stops renewing; its lease lapses and the demand path collapses inward toward $\ell(\contract)$. A subscriber whose lease expires must re-subscribe. This interest-gating keeps the live subscription set proportional to active demand rather than to accumulated cache: orphaned paths die out within a small multiple of the lease period.

\subsection{Relationship to Other Mechanisms}

\paragraph{Versus operation-based CRDTs.} An operation-based CRDT propagates individual operations and relies on a causal-broadcast layer to ensure exactly-once, causally ordered delivery. The protocol here propagates state and tolerates arbitrary loss, reordering, and duplication; the merge function does the work that causal broadcast would otherwise have to.

\paragraph{Versus Merkle anti-entropy.} Distributed key-value stores commonly use Merkle-tree anti-entropy \citep{decandia2007dynamo} to discover differences between replicas. That is a special case of the protocol here in which the summary is a hash tree and the delta is the subtree the other replica is missing. A contract whose data structure is amenable to a more compact summary (a list of identifiers, a counter value) pays a single round trip rather than $O(\log n)$ of them.

\paragraph{Versus rsync.} \texttt{rsync} \citep{tridgell1999rsync} reconciles files by exchanging block-level checksums. It is generic but oblivious to application structure. Summary/delta is the same shape of protocol lifted to a level where the application defines the granularity.

\paragraph{Versus delta-state CRDTs.} The closest formal analogue is the delta-state CRDT literature \citep{almeida2018delta}, which ships only changes between replicas of a fixed-lattice CRDT. The differences here are that the lattice and the delta-encoding are both supplied by the contract, and the result is embedded in a small-world routing layer rather than a fully-connected replica set.
